\chapter{Examples and counterexamples}
\label{chap:examples}

In this chapter we test the general notions of the previous chapters on concrete examples. The
positive examples (closed subspaces, open subsets, graphs, the unit sphere of an inner product
space, the unitary group of a Banach $*$-algebra) show that embedded and split submanifolds,
fields of projectors and their geodesics behave as expected; in all of them the slice charts are
explicit, so that no inverse function theorem is needed. The counterexamples separate the
notions of submanifold and immersion (Definitions~\ref{def:weak-immersion},
\ref{def:immersion}, \ref{def:split-immersion} and~\ref{def:initial-submanifold}): a weak
immersion that is not an immersion ($\ell^1 \to \ell^2$), an injective immersion that is not
initial (the figure eight), and an immersion between Banach spaces that is not split
($c_0 \subseteq \ell^\infty$, by Phillips' lemma). The last section prepares the ground for
groups of diffeomorphisms: we construct the Fréchet space $C^\infty_{per}$ of smooth periodic
functions and identify the universal cover of $\Diff^+(S^1)$ with an open convex subset of it.

Throughout, we keep track of completeness. A model space is a real locally convex Hausdorff
space, not necessarily complete; when an example is a Banach, Hilbert or Fréchet manifold we say
so, and we say where completeness is used.

\section{Closed subspaces, open subsets and graphs}

The simplest submanifolds are linear. Let $F$ be a locally convex space and $S \subseteq F$ a
linear subspace.

\begin{proposition}[Closed subspaces]
  \label{prop:ex-subspace}
  \lean{LCManifolds.isSliceChart_refl, LCManifolds.isEmbeddedSubmanifold_subspace,
    LCManifolds.isSplitSubmanifold_subspace, LCManifolds.tangentSpace_subspace}
  \leanok
  \uses{def:embedded-submanifold, def:split-submanifold, def:tangent-space-submanifold}
  If $S$ is closed, then $S$ is an embedded submanifold of $F$ modeled on $S$, with the identity
  of $F$ as slice chart at every point, and $T_pS = S$ for every $p \in S$. If $S$ is
  complemented, $S$ is a split submanifold.
\end{proposition}

\begin{proof}
  \leanok
  The identity belongs to the maximal atlas of $F$, and $q \in S$ if and only if
  $\id(q) \in S$. The differential of the identity is the identity, so the tangent space read
  in this chart is $S$.
\end{proof}

\begin{proposition}[Linear inclusions]
  \label{prop:ex-subspace-immersion}
  \lean{LCManifolds.isImmersion_subtypeL, LCManifolds.isSplitImmersion_subtypeL_iff}
  \leanok
  \uses{def:immersion, def:split-immersion}
  If $S$ is closed, the inclusion $\iota \colon S \to F$ is an immersion. It is a split
  immersion if and only if $S$ is complemented in $F$.
\end{proposition}

\begin{proof}
  \leanok
  In the identity charts the differential of $\iota$ at every point is $\iota$ itself, which is
  a topological isomorphism onto the closed subspace $S$. A continuous left inverse of $\iota$
  is the same thing as a continuous projection of $F$ onto $S$.
\end{proof}

A continuous projection $P$ of $F$ onto $S$ ($P^2 = P$, $\operatorname{ran} P = S$) gives the
simplest field of projectors: the constant one.

\begin{proposition}[The flat connection of a subspace]
  \label{prop:ex-subspace-flat}
  \lean{LCManifolds.projectorField_subspace, LCManifolds.symb_subspace_refl, LCManifolds.isGeodesic_subspace_iff}
  \leanok
  \uses{prop:ex-subspace, def:projector-field, thm:projector-connection, def:geodesic}
  Let $S$ be closed and $P$ a continuous projection onto $S$. Then $E_p = P$ for all $p \in S$
  is a field of projectors, the symbols of its connection vanish in the identity chart, and a
  $\Ck{2}$ curve $\alpha$ in $S$ is a geodesic if and only if $\ddot\alpha(t) = 0$ for all $t$.
\end{proposition}

\begin{proof}
  \leanok
  \uses{thm:projector-geodesics}
  Smoothness of $(p, w) \mapsto P w$ is that of a continuous linear map. The symbols involve the
  second differential of the inclusion, which is zero. For the geodesics, the general criterion
  says that $\alpha$ is a geodesic iff $P\ddot\alpha = 0$; since $\ddot\alpha$ takes values in
  $S$ and $P$ is the identity on $S$, this is $\ddot\alpha = 0$.
\end{proof}

So the connection is flat and its geodesics are the curves of zero acceleration, that is, the
affine lines $t \mapsto x + t v$ (this last reformulation is the usual one and is not stated
separately in Lean).

\begin{proposition}[Open subsets]
  \label{prop:ex-open}
  \lean{LCManifolds.isEmbeddedSubmanifold_open}
  \leanok
  \uses{def:embedded-submanifold, def:mb-manifold}
  An open subset $U$ of a manifold $M$ modeled on $E$ is an embedded submanifold of $M$, modeled
  on the whole space $E$.
\end{proposition}

\begin{proof}
  \leanok
  The restriction of the chart at $p$ to $U$ is again a chart (the Michal--Bastiani groupoid is
  closed under restriction to open sets), and every point of its domain lies in $U$.
\end{proof}

Let now $E, G$ be locally convex spaces, $U \subseteq E$ open and $g \colon U \to G$ smooth.
The graph of $g$ is
\[
  \Gamma_g = \{(x, y) \in E \times G : x \in U,\ y = g(x)\},
\]
and its model is the horizontal subspace $E \times 0$, the kernel of the second projection.

\begin{proposition}[Graphs]
  \label{prop:ex-graph}
  \lean{LCManifolds.graph, LCManifolds.horizontal, LCManifolds.closedComplemented_horizontal, LCManifolds.graphChart, LCManifolds.graphChart_mem, LCManifolds.isSplitSubmanifold_graph}
  \leanok
  \uses{def:split-submanifold, def:mb-groupoid}
  The map $\Phi(x, y) = (x, y - g(x))$, defined on $U \times G$, is a chart of $E \times G$ with
  inverse $\Phi^{-1}(x, z) = (x, z + g(x))$, and $(x, y) \in \Gamma_g$ if and only if
  $\Phi(x, y) \in E \times 0$. Hence $\Gamma_g$ is a split submanifold of $E \times G$ modeled on
  $E \times 0$.
\end{proposition}

\begin{proof}
  \leanok
  Both $\Phi$ and $\Phi^{-1}$ are built from $g$, projections and sums, hence they are smooth on
  $U \times G$. The subspace $E \times 0$ is closed and complemented by $0 \times G$, with
  projection $(x, y) \mapsto (x, 0)$.
\end{proof}

\begin{remark}
  In finite dimensions one usually proves that graphs (and level sets of submersions) are
  submanifolds through the inverse function theorem. That theorem is not available for
  Michal--Bastiani maps between locally convex spaces, and it is not needed here: the chart
  $\Phi$ and its inverse are written down explicitly. The same strategy is used for the sphere
  and the unitary group below.
\end{remark}

\section{The unit sphere of an inner product space}

Let $H$ be a real inner product space and $e \in H$ a unit vector. We write
$e^\perp = \{w : \langle w, e\rangle = 0\}$, a closed subspace of $H$, and
$S = \{q \in H : \|q\| = 1\}$.

\begin{remark}[Completeness]
  We do not assume that $H$ is complete. If $H$ is complete, $H$ is a Hilbert manifold and $S$ is
  a Hilbert submanifold. If $H$ is not complete, $H$ is still a manifold in our sense (model
  spaces need not be complete), but neither $H$ nor $S$ is a Hilbert or Banach manifold. The
  construction below does not use completeness, so it covers both cases; the statements about
  Hilbert manifolds are the case of a complete $H$.
\end{remark}

On the open hemisphere around $e$ we use the explicit chart
\[
  \varphi_e(q) = q - \sqrt{1 - \|q - \langle q, e\rangle e\|^2}\; e, \qquad
  \varphi_e^{-1}(r) = r + \sqrt{1 - \|r - \langle r, e\rangle e\|^2}\; e,
\]
defined on $\{q : \|q - \langle q, e\rangle e\| < 1,\ \langle q, e\rangle > 0\}$. By Pythagoras,
$\|q\|^2 = \|q - \langle q, e\rangle e\|^2 + \langle q, e\rangle^2$, so on this set $\|q\| = 1$
exactly when the height $\sqrt{1 - \|q - \langle q,e\rangle e\|^2}$ equals $\langle q, e\rangle$,
that is, when $\varphi_e(q) \in e^\perp$. At another point $p \in S$ we first apply the
Householder reflection $R_p$ exchanging $p$ and $e$ (the reflection across
$(p - e)^\perp$), a linear isometry, and then $\varphi_e$.

\begin{lemma}[The charts of the sphere]
  \label{lem:ex-sphere-charts}
  \lean{LCManifolds.Sphere.chart, LCManifolds.Sphere.chart_mem, LCManifolds.Sphere.norm_eq_one_iff, LCManifolds.Sphere.refl, LCManifolds.Sphere.chartAt', LCManifolds.Sphere.chartAt'_mem, LCManifolds.Sphere.mem_chartAt'_source}
  \leanok
  \uses{def:mb-groupoid}
  The map $\varphi_e$ is a smooth chart of $H$ with smooth inverse; for $q$ in its domain,
  $\|q\| = 1$ if and only if $\varphi_e(q) \in e^\perp$. For every unit vector $p$, the chart
  $\varphi_e \circ R_p$ is smooth and its domain contains $p$.
\end{lemma}

\begin{proof}
  \leanok
  The square root is smooth where its argument is positive, and the remaining operations are
  continuous linear or bilinear; Fréchet smoothness implies Michal--Bastiani smoothness. The
  identity $\varphi_e^{-1}\varphi_e = \id$ holds because the component of $q$ orthogonal to $e$,
  and hence the height, does not change when we add a multiple of $e$. The equivalence is
  Pythagoras, and $R_p p = e$ lies in the domain of $\varphi_e$.
\end{proof}

\begin{theorem}[The sphere is an embedded submanifold]
  \label{thm:ex-sphere}
  \lean{LCManifolds.isEmbeddedSubmanifold_sphere, LCManifolds.tangentSpace_sphere}
  \leanok
  \uses{lem:ex-sphere-charts, def:embedded-submanifold, def:tangent-space-submanifold}
  The unit sphere $S$ is an embedded submanifold of $H$ modeled on $e^\perp$, and
  $T_pS = p^\perp$ for every $p \in S$.
\end{theorem}

\begin{proof}
  \leanok
  Since $R_p$ is an isometry, $\|q\| = 1$ iff $\|R_p q\| = 1$ iff
  $\varphi_e(R_p q) \in e^\perp$. At $e$ the height has zero derivative, so
  $\dd\varphi_e(e) = \id$, and by the chain rule $\dd(\varphi_e \circ R_p)(p) = R_p$. Then
  $w \in T_pS$ iff $\langle e, R_p w\rangle = 0$ iff $\langle R_p p, R_p w\rangle = 0$ iff
  $\langle p, w\rangle = 0$.
\end{proof}

The orthogonal projections onto the tangent spaces give a field of projectors, and the geodesics
of its connection are the curves whose acceleration is normal to the sphere.

\begin{theorem}[Geodesics of the sphere]
  \label{thm:ex-sphere-geodesics}
  \lean{LCManifolds.sphereProj, LCManifolds.projectorField_sphere, LCManifolds.isGeodesic_sphere_iff, LCManifolds.greatCircle, LCManifolds.isGeodesic_greatCircle}
  \leanok
  \uses{thm:ex-sphere, def:projector-field, thm:projector-connection, def:geodesic}
  The maps $E_p w = w - \langle w, p\rangle p$, $p \in S$, form a field of projectors onto
  the tangent spaces of $S$. A curve $\alpha$ in $S$ whose ambient curve is smooth is a geodesic
  of the associated connection if and only if
  \[
    \ddot\alpha(t) = \langle \ddot\alpha(t), \alpha(t)\rangle\, \alpha(t) \quad
    \text{for all } t .
  \]
  In particular, if $p, v \in H$ are unit vectors with $\langle p, v\rangle = 0$, the great
  circle $t \mapsto \cos t\, p + \sin t\, v$ is a geodesic.
\end{theorem}

\begin{proof}
  \leanok
  \uses{thm:projector-geodesics}
  It is immediate that $E_p$ is linear, idempotent, with range $p^\perp$. In a slice chart,
  $(x, w) \mapsto E_{\psi(x)} w$ is smooth because the parametrization $\psi$ is. By the general
  criterion, $\alpha$ is a geodesic iff $E_{\alpha(t)}\ddot\alpha(t) = 0$, which is the stated
  equation. For the great circle, $\ddot\alpha = -\alpha$ and $\|\alpha\| = 1$.
\end{proof}

\section{Weak immersions, injective immersions and split immersions}

The notions of immersion form a chain: split immersions are immersions,
immersions are weak immersions, and injective immersions need not be initial. In finite
dimensions the first two notions coincide. The three examples of this section show that every
implication is strict, and that the first two separations occur already between Banach spaces.

\subsection{A weak immersion that is not an immersion}

Let $\iota \colon \ell^1 \to \ell^2$ be the inclusion of real sequences. Both spaces are Banach
spaces, hence manifolds with a single chart.

\begin{proposition}[$\ell^1 \to \ell^2$]
  \label{prop:ex-l1-l2}
  \lean{LCManifolds.lpInclusion, LCManifolds.lpInclusion_injective, LCManifolds.not_boundedBelow_lpInclusion, LCManifolds.isWeakImmersion_lpInclusion, LCManifolds.not_isImmersion_lpInclusion}
  \leanok
  \uses{def:weak-immersion, def:immersion}
  The inclusion $\iota$ is continuous, linear and injective, hence a weak immersion. It is not
  bounded below, hence it is not an immersion.
\end{proposition}

\begin{proof}
  \leanok
  Continuity follows from the closed graph theorem: the graph of $\iota$ is closed because the
  coordinate evaluations are continuous on both spaces. The differential of $\iota$ at every
  point is $\iota$, which is injective. For $x_n = e_0 + \dots + e_{n-1}$ we have
  $\|x_n\|_1 = n$ and $\|x_n\|_2 = \sqrt n$, so no inequality $m\|x\|_1 \le \|\iota x\|_2$ with
  $m > 0$ can hold, and $\iota$ is not a topological isomorphism onto its range.
\end{proof}

\begin{remark}
  The range of $\iota$ is dense in $\ell^2$ and not closed. In finite dimensions an injective
  linear map is automatically a topological isomorphism onto a closed subspace; in infinite
  dimensions the condition that the differential be an embedding with closed range is a genuine
  additional requirement, and this is the reason to distinguish weak immersions from
  immersions.
\end{remark}

\subsection{An injective immersion that is not initial: the figure eight}

Let $\theta(t) = 2\arctan t \in (-\pi, \pi)$ and
\[
  j \colon \R \to \R^2, \qquad j(t) = (\sin\theta(t), \sin 2\theta(t)).
\]
As $t$ runs over $\R$, $j(t)$ traces the figure eight (a lemniscate), passing through the origin
at $t = 0$ and tending to the origin again as $t \to \pm\infty$.

\begin{proposition}[The figure eight is immersed]
  \label{prop:ex-lemniscate-immersed}
  \lean{LCManifolds.Lemniscate.j, LCManifolds.Lemniscate.j_injective, LCManifolds.Lemniscate.vel_ne_zero, LCManifolds.Lemniscate.isImmersedSubmanifold_lemniscate}
  \leanok
  \uses{def:immersion}
  The map $j$ is smooth, injective and has nonvanishing velocity; hence it is an injective
  immersion.
\end{proposition}

\begin{proof}
  \leanok
  On $(-\pi, \pi)$ the pair $(\sin a, \sin 2a)$ determines $a$, which gives injectivity. If
  $\dot j(t) = 0$, then $\cos\theta = 0$ and $\cos 2\theta = 0$ at the same time, which is
  impossible. A nonzero vector spans a closed line, so $v \mapsto v\,\dot j(t)$ is a
  topological isomorphism onto a closed subspace.
\end{proof}

\begin{theorem}[The figure eight is not initial]
  \label{thm:ex-lemniscate-not-initial}
  \lean{LCManifolds.Lemniscate.f, LCManifolds.Lemniscate.g, LCManifolds.Lemniscate.j_g, LCManifolds.Lemniscate.one_lt_g, LCManifolds.Lemniscate.g_zero, LCManifolds.Lemniscate.not_isInitialSubmanifold_lemniscate}
  \leanok
  \uses{prop:ex-lemniscate-immersed, def:initial-submanifold}
  Let $\sigma(s) = \tfrac12\arctan s$ and $f(s) = (\sin\sigma, -\sin 2\sigma)$. Then $f$ is a
  smooth curve in $\R^2$ with values in $j(\R)$, but its unique lift $g$ with $j \circ g = f$ is
  not continuous at $0$. Hence $j$ is not an initial submanifold.
\end{theorem}

\begin{proof}
  \leanok
  The curve $f$ crosses the origin along the branch that $j$ reaches only as $t \to \pm\infty$.
  Explicitly, $g = \tan(\lambda/2)$ with $\lambda(s) = \pi - \sigma(s)$ for $s > 0$,
  $\lambda(s) = -\pi - \sigma(s)$ for $s < 0$ and $\lambda(0) = 0$; trigonometric identities give
  $j \circ g = f$. For $s > 0$ we have $g(s) > 1$, while $g(0) = 0$. If $j$ were initial, $g$
  would be smooth, in particular continuous.
\end{proof}

\begin{remark}
  This is the classical finite dimensional example; it shows that the condition defining
  initial submanifolds is not a consequence of being an injective immersion, and has nothing to
  do with infinite dimensions. The topology of $\R$ transported by $j$ is strictly finer than the
  topology induced by $\R^2$ on the figure eight.
\end{remark}

\subsection{An immersion that is not split: \texorpdfstring{$c_0 \subseteq \ell^\infty$}{c0 in l-infinity}}

We realize $\ell^\infty$ as the Banach space of bounded real sequences and $c_0$ as the closed
subspace of sequences tending to $0$. The classical obstruction is the following result of
Phillips \cite{Phi40}; we formalize the short proof of Whitley \cite{Whi66}.

\begin{theorem}[Phillips' lemma]
  \label{thm:ex-phillips}
  \lean{LCManifolds.Phillips.Linf, LCManifolds.Phillips.c0, LCManifolds.Phillips.isClosed_c0, LCManifolds.Phillips.not_closedComplemented_c0}
  \leanok
  There is no continuous linear projection of $\ell^\infty$ onto $c_0$.
\end{theorem}

The proof rests on an uncountable family of almost disjoint subsets of $\N$, and on the fact
that a functional vanishing on $c_0$ can only see countably many of them.

\begin{lemma}[An almost disjoint family]
  \label{lem:ex-almost-disjoint}
  \lean{LCManifolds.Phillips.N, LCManifolds.Phillips.infinite_N, LCManifolds.Phillips.finite_N_inter, LCManifolds.Phillips.ind, LCManifolds.Phillips.ind_not_mem_c0}
  \leanok
  For each $A \subseteq \N$ there is an infinite set $N_A \subseteq \N$ such that
  $N_A \cap N_B$ is finite whenever $A \ne B$. The indicator $\chi_A \in \ell^\infty$ of $N_A$
  does not belong to $c_0$.
\end{lemma}

\begin{proof}
  \leanok
  Index the finite binary words by $\N$, and let $N_A$ be the set of indices of the initial
  segments of the indicator function of $A$. If $A$ and $B$ differ at $i_0$, their initial
  segments of length larger than $i_0$ are different, so $N_A \cap N_B$ is finite. Since $N_A$
  is infinite, $\chi_A$ takes the value $1$ infinitely often.
\end{proof}

\begin{lemma}[Whitley]
  \label{lem:ex-whitley}
  \lean{LCManifolds.Phillips.abs_le_of_eventually, LCManifolds.Phillips.countable_ne_zero}
  \leanok
  \uses{lem:ex-almost-disjoint}
  Let $f$ be a continuous linear functional on $\ell^\infty$ that vanishes on $c_0$. If
  $|x_k| \le 1$ for all but finitely many $k$, then $|f(x)| \le \|f\|$. Consequently,
  $f(\chi_A) \ne 0$ only for countably many $A \subseteq \N$.
\end{lemma}

\begin{proof}
  \leanok
  For the first claim, modify $x$ on a finite set to obtain $y$ with $\|y\| \le 1$; then
  $x - y \in c_0$ and $f(x) = f(y)$. For the second, suppose $|f(\chi_A)| > 1/m$ for every $A$ in
  a finite set $\mathcal F$, and let $x = \sum_{A \in \mathcal F} \pm\chi_A$, with the signs of
  $f(\chi_A)$. Off the finite set where two of the $N_A$ meet, $|x_k| \le 1$; hence
  $\#\mathcal F/m < f(x) \le \|f\|$. So each set $\{A : |f(\chi_A)| > 1/m\}$ is finite.
\end{proof}

\begin{proof}[Proof of Theorem~\ref{thm:ex-phillips}]
  \leanok
  \uses{lem:ex-whitley}
  Suppose $P$ is a continuous projection onto $c_0$, and apply Lemma~\ref{lem:ex-whitley} to the
  functionals $f_n = e_n \circ (1 - P)$, where $e_n$ is the $n$-th coordinate. The subsets of
  $\N$ are uncountable, so some $A$ avoids the countable union of the exceptional sets. Then
  $(1 - P)\chi_A = 0$, that is, $\chi_A = P\chi_A \in c_0$, a contradiction.
\end{proof}

\begin{corollary}[Embedded but not split]
  \label{cor:ex-c0}
  \lean{LCManifolds.Phillips.isEmbeddedSubmanifold_c0, LCManifolds.Phillips.not_isSplitSubmanifold_c0, LCManifolds.Phillips.isImmersion_c0, LCManifolds.Phillips.not_isSplitImmersion_c0}
  \leanok
  \uses{thm:ex-phillips, prop:ex-subspace, prop:ex-subspace-immersion,
    def:embedded-submanifold, def:split-submanifold, def:immersion, def:split-immersion}
  The subspace $c_0$ is an embedded submanifold of $\ell^\infty$ that is not split, and the
  inclusion $c_0 \to \ell^\infty$ is an immersion that is not a split immersion.
\end{corollary}

\begin{proof}
  \leanok
  Combine Propositions~\ref{prop:ex-subspace} and~\ref{prop:ex-subspace-immersion} with
  Theorem~\ref{thm:ex-phillips}.
\end{proof}

\begin{remark}
  In Hilbert spaces every closed subspace is complemented, so the distinction between embedded
  and split submanifolds, and between immersions and split immersions, is invisible there. For
  Banach spaces it is not, and Corollary~\ref{cor:ex-c0} gives the standard witness.
\end{remark}

\section{The unitary group of a Banach \texorpdfstring{$*$}{*}-algebra}

Let $A$ be a real Banach algebra with a continuous involution $a \mapsto a^*$ (in Lean:
a complete normed $\R$-algebra which is a $*$-ring and a $*$-module, with continuous $*$). For
instance $A = B(H)$, the bounded operators on a complex Hilbert space $H$, seen as a real
algebra. The unitary group is
\[
  U(A) = \{u \in A : u^* u = 1 = u u^*\},
\]
and its model is the space $\mathfrak u(A) = \{X \in A : X^* = -X\}$ of skew-adjoint elements.

\begin{remark}[Completeness]
  Here $A$ is complete, so $A$ is a Banach manifold and $U(A)$ is a Banach submanifold of it.
  Completeness is used twice: the invertible elements form an open set, and inversion is smooth
  on it. Both facts may fail in an incomplete normed algebra.
\end{remark}

The slice charts are Cayley transforms. On the open set $\{a : a + 1 \text{ invertible}\}$ we
let
\[
  \varphi(a) = 1 - 2(a + 1)^{-1} = (a - 1)(a + 1)^{-1}, \qquad
  \varphi^{-1}(X) = 2(1 - X)^{-1} - 1 = (1 + X)(1 - X)^{-1}.
\]
At other points $p \in U(A)$ we compose with the left translation $a \mapsto p^* a$, a linear
homeomorphism of $A$ with inverse $a \mapsto pa$.

\begin{lemma}[The Cayley transform]
  \label{lem:ex-cayley}
  \lean{LCManifolds.Cayley.cay, LCManifolds.Cayley.cayInv, LCManifolds.Cayley.chart, LCManifolds.Cayley.chart_mem, LCManifolds.Cayley.mem_unitary_iff, LCManifolds.Cayley.chartAt', LCManifolds.Cayley.isSliceChart}
  \leanok
  \uses{def:mb-groupoid}
  The Cayley transform $\varphi$ is a smooth chart of $A$ with smooth inverse. If $a + 1$ is
  invertible, then $a \in U(A)$ if and only if $\varphi(a) \in \mathfrak u(A)$. For every
  $p \in U(A)$, the chart $a \mapsto \varphi(p^* a)$ is a slice chart for $U(A)$ around $p$.
\end{lemma}

\begin{proof}
  \leanok
  The inverse of $\varphi$ is checked by a direct computation with units, and smoothness follows
  from the smoothness of inversion on the open set of units. The equivalence is pure algebra:
  writing $s = a + 1$ and using $(s^{-1})^* = (s^*)^{-1}$, the condition
  $\varphi(a)^* = -\varphi(a)$ unwinds to $a^* a = 1 = a a^*$; no spectral theory is needed.
  Finally $p^* p + 1 = 2$ is invertible, so $p$ lies in the domain of the translated chart.
\end{proof}

\begin{theorem}[The unitary group is a split submanifold]
  \label{thm:ex-unitary}
  \lean{LCManifolds.closedComplemented_skewAdjoint, LCManifolds.isSplitSubmanifold_unitary, LCManifolds.tangentSpace_unitary, LCManifolds.isSplitSubmanifold_unitary_clm}
  \leanok
  \uses{lem:ex-cayley, def:split-submanifold, def:tangent-space-submanifold}
  The unitary group $U(A)$ is a split submanifold of $A$ modeled on $\mathfrak u(A)$, and
  \[
    T_pU(A) = \{w \in A : p^* w \in \mathfrak u(A)\} = p\,\mathfrak u(A) .
  \]
  In particular this holds for the unitary group of $B(H)$, $H$ a complex Hilbert space.
\end{theorem}

\begin{proof}
  \leanok
  The skew-adjoint elements are closed and complemented by the self-adjoint ones, with
  projection $X \mapsto (X - X^*)/2$. The derivative of $\varphi$ at $1$ is $w \mapsto w/2$, so
  the derivative of the translated chart at $p$ is $w \mapsto p^* w/2$, and the tangent space
  is read off as in Theorem~\ref{thm:ex-sphere}.
\end{proof}

The natural projector onto $T_pU(A)$ is $E_p w = \tfrac12(w - p w^* p)$; writing $X = p^* w$, it
is $E_p w = p\,\tfrac12(X - X^*)$, the skew-adjoint part translated back to $p$.

\begin{theorem}[Geodesics of the unitary group]
  \label{thm:ex-unitary-geodesics}
  \lean{LCManifolds.unitaryProj, LCManifolds.unitaryProj_idem, LCManifolds.projectorField_unitary, LCManifolds.isGeodesic_unitary_iff, LCManifolds.expCurve, LCManifolds.isGeodesic_expCurve}
  \leanok
  \uses{thm:ex-unitary, def:projector-field, thm:projector-connection, def:geodesic}
  The maps $E_p w = \tfrac12(w - p w^* p)$, $p \in U(A)$, form a field of projectors onto the
  tangent spaces of $U(A)$. A curve $\alpha$ in $U(A)$ whose ambient curve is smooth is a
  geodesic of the associated connection if and only if
  \[
    \ddot\alpha(t) = \alpha(t)\, \ddot\alpha(t)^*\, \alpha(t) \quad \text{for all } t .
  \]
  For $p \in U(A)$ and $X \in \mathfrak u(A)$, the curve $t \mapsto p\, e^{tX}$ is a geodesic.
\end{theorem}

\begin{proof}
  \leanok
  \uses{thm:projector-geodesics}
  Using $p^*p = 1 = pp^*$ one checks $p (E_p w)^* p = \tfrac12(p w^* p - w)$, whence
  $E_p^2 = E_p$, and $p^* E_p w = \tfrac12(p^* w - w^* p)$ is skew-adjoint. The geodesic
  equation is the general criterion $E_{\alpha}\ddot\alpha = 0$. For $\alpha(t) = p e^{tX}$,
  $e^{tX}$ is unitary and commutes with $X$, and $\ddot\alpha = p e^{tX}X^2$; a short
  computation gives $\alpha\ddot\alpha^*\alpha = p e^{tX} X^2$.
\end{proof}

\begin{remark}
  The structure of Lie group of $U(A)$ (smoothness of the product and of the inversion in these
  charts) is not addressed in this chapter; we only use $U(A)$ as a submanifold of $A$. The
  connection of Theorem~\ref{thm:ex-unitary-geodesics} is defined directly from the field of
  projectors, with no metric behind it. The header of the Lean file also mentions $B(H)$ for a
  real Hilbert space $H$; the formalized statement \texttt{isSplitSubmanifold\_unitary\_clm}
  covers complex Hilbert spaces.
\end{remark}

\section{Smooth periodic functions and \texorpdfstring{$\Diff^+(S^1)$}{Diff+(S1)}}

\subsection{The Fréchet space \texorpdfstring{$C^\infty_{per}$}{C-infinity-per}}

Let $C^\infty_{per}$ be the space of smooth $1$-periodic functions $u \colon \R \to \R$ (smooth
functions on the circle $S^1 = \R/\mathbb Z$), with the topology of uniform convergence of all
the derivatives, that is, the topology of the seminorms
\[
  p_k(u) = \sup_{x \in \R} |u^{(k)}(x)|, \qquad k \in \N .
\]

\begin{definition}[$C^\infty_{per}$]
  \label{def:ex-cinfper}
  \lean{LCManifolds.CInfPer, LCManifolds.CInfPer.J, LCManifolds.CInfPer.isEmbedding_J, LCManifolds.CInfPer.seminorm}
  \leanok
  The topology and the uniform structure of $C^\infty_{per}$ are those induced by the injective
  linear map
  \[
    J \colon C^\infty_{per} \to \prod_{k \in \N} BC(\R), \qquad J(u) = (u, u', u'', \dots),
  \]
  into a countable product of Banach spaces of bounded continuous functions. The seminorm $p_k$
  is the sup norm of the $k$-th component of $J$.
\end{definition}

All the structure is inherited from the product: $C^\infty_{per}$ is a Hausdorff locally convex
topological vector space, and it is metrizable, since the product is countable. In Lean these
are instances on \texttt{CInfPer}.

\begin{proposition}[Convergence and seminorms]
  \label{prop:ex-cinfper-tendsto}
  \lean{LCManifolds.CInfPer.tendsto_iff, LCManifolds.CInfPer.norm_iteratedDeriv_le_seminorm, LCManifolds.CInfPer.seminorm_le_iff}
  \leanok
  \uses{def:ex-cinfper}
  A net $u_i$ converges to $u$ in $C^\infty_{per}$ if and only if $u_i^{(k)} \to u^{(k)}$
  uniformly on $\R$ for every $k$. Moreover $p_k(u) \le C$ if and only if
  $|u^{(k)}(x)| \le C$ for all $x$.
\end{proposition}

\begin{proof}
  \leanok
  $J$ is a uniform embedding and convergence in a product is convergence of every component.
\end{proof}

\begin{theorem}[$C^\infty_{per}$ is a Fréchet space]
  \label{thm:ex-cinfper-frechet}
  \lean{LCManifolds.CInfPer.isClosed_range_J}
  \leanok
  \uses{def:ex-cinfper}
  The range of $J$ is closed. Consequently $C^\infty_{per}$ is complete, and it is a Fréchet
  space: locally convex, Hausdorff, metrizable and complete.
\end{theorem}

\begin{proof}
  \leanok
  If $u_n^{(k)} \to g_k$ uniformly for every $k$, the theorem on uniform limits of derivatives
  shows that $g_0$ is differentiable with $g_0' = g_1$, and inductively $g_0^{(k)} = g_k$; so
  $g_0$ is smooth, periodic, and $J(g_0) = (g_k)_k$. A closed subspace of the complete space
  $\prod_k BC(\R)$ is complete, and $J$ is a uniform embedding.
\end{proof}

\begin{remark}
  The space $C^\infty_{per}$ is not normable, so it is a Fréchet space that is not a Banach
  space. Non-normability is the standard fact that a normable Fréchet space has a bounded
  neighbourhood of $0$; it is not formalized here. Completeness, which we do prove, is what
  makes the manifolds modeled on $C^\infty_{per}$ below Fréchet manifolds in the strict sense.
\end{remark}

\subsection{The universal cover of \texorpdfstring{$\Diff^+(S^1)$}{Diff+(S1)}}

Let $\pi \colon \R \to S^1 = \R/\mathbb Z$ be the projection. We recall the classical facts
that relate circle maps to their lifts; they are \emph{not} formalized, and the formalization
works directly with the lifts.

\begin{remark}[From the circle to its lifts]
  \label{rem:ex-covering}
  The map $\pi$ is a covering map and a local diffeomorphism.
  \begin{itemize}
    \item \emph{Existence and uniqueness of lifts.} Since $\R$ is simply connected, every
    continuous $f \colon S^1 \to S^1$ has a continuous lift $\varphi \colon \R \to \R$,
    $\pi \circ \varphi = f \circ \pi$, by the lifting criterion \cite[\S1.3, Prop.~1.33]{Hat02}
    (see also the homotopy lifting property, \cite[Prop.~1.30]{Hat02}), and two lifts differ by
    an integer constant: integer translations are deck transformations, and a lift is
    determined by its value at one point \cite[Prop.~1.34]{Hat02}.
    \item \emph{Degree.} $\varphi(x + 1) - \varphi(x)$ is a continuous integer valued function,
    hence a constant $d$, the degree of $f$ \cite{KH95}. A homeomorphism has degree $\pm1$, and it
    preserves the orientation exactly when $d = 1$, i.e.\ $\varphi(x + 1) = \varphi(x) + 1$;
    then $\varphi = \id + u$ with $u$ $1$-periodic.
    \item \emph{Smoothness.} Since $\pi$ is a local diffeomorphism, $\varphi$ is smooth if and
    only if $f$ is, and $\varphi$ is a diffeomorphism of $\R$ exactly when $f$ is a
    diffeomorphism of $S^1$ (the inverse of $f$ lifts to the inverse of $\varphi$).
  \end{itemize}
  So $\Diff^+(S^1) \cong \Omega/\mathbb Z$, where $n \in \mathbb Z$ acts by $u \mapsto u + n$ and
  $\Omega$ is the set defined below. This is the description of Herman
  \cite[Ch.~I, 1.7]{Her79}, where $D^r(T^1) = \{\id + \varphi\}$ is called the universal cover
  of $\Diff^r_+(T^1)$; Hamilton \cite[Part~I, Example~4.2.6]{Ham82} parametrizes
  $\Diff^+(S^1)$ instead by $S^1 \times \{f' > 0,\ \int f' = 2\pi\}$. Since $\Omega$ is convex,
  hence contractible, $\Omega \to \Diff^+(S^1)$ is the universal cover
  $\widetilde{\Diff}^+(S^1)$ \cite[Introduction (Remarque importante)]{Her79}. For Lie groups
  of diffeomorphisms in the locally convex setting see \cite{GN26}. In Mathlib, the monotone
  lifts with $\varphi(x+1) = \varphi(x) + 1$ of circle maps appear as
  \texttt{CircleDeg1Lift}.
\end{remark}

An orientation preserving diffeomorphism of the circle thus lifts to a diffeomorphism
$\varphi = \id + u$ of $\R$ with $u \in C^\infty_{per}$ and $\varphi' = 1 + u' > 0$.

\begin{definition}[The domain $\Omega$]
  \label{def:ex-omega}
  \lean{LCManifolds.DiffCircle.domain, LCManifolds.DiffCircle.Ω, LCManifolds.DiffCircle.lift}
  \leanok
  \uses{def:ex-cinfper}
  We let
  \[
    \Omega = \{u \in C^\infty_{per} : 1 + u'(x) > 0 \text{ for all } x \in \R\},
  \]
  and for $u \in C^\infty_{per}$ we let $\varphi_u = \id + u$ be its lift.
\end{definition}

\begin{proposition}[$\Omega$ is open and convex]
  \label{prop:ex-omega-open}
  \lean{LCManifolds.DiffCircle.exists_pos_le, LCManifolds.DiffCircle.isOpen_domain, LCManifolds.DiffCircle.convex_domain, LCManifolds.DiffCircle.zero_mem}
  \leanok
  \uses{def:ex-omega, thm:ex-cinfper-frechet, prop:ex-open, def:mb-manifold}
  The set $\Omega$ is open and convex in $C^\infty_{per}$, and contains $0$. Hence $\Omega$ is a
  manifold modeled on $C^\infty_{per}$; since the model is complete, it is a Fréchet manifold.
  It is not a Banach manifold, since the model is not normable (an unformalized fact).
\end{proposition}

\begin{proof}
  \leanok
  For $u \in \Omega$, the continuous periodic function $1 + u'$ attains a positive minimum $c$ on
  the compact circle. If $p_1(v - u) < c$, then $1 + v' > 0$, so $\Omega$ contains a
  neighbourhood of $u$. Convexity is the identity
  $1 + (au + bv)' = a(1 + u') + b(1 + v')$ for $a, b \ge 0$, $a + b = 1$. An open subset of a
  locally convex space is a manifold with the restricted identity chart.
\end{proof}

\begin{proposition}[Lifts are diffeomorphisms]
  \label{prop:ex-lift-diffeo}
  \lean{LCManifolds.DiffCircle.lift_add_one, LCManifolds.DiffCircle.strictMono_lift, LCManifolds.DiffCircle.surjective_lift, LCManifolds.DiffCircle.liftHomeomorph, LCManifolds.DiffCircle.contDiff_liftHomeomorph_symm, LCManifolds.DiffCircle.liftHomeomorph_symm_add_one}
  \leanok
  \uses{def:ex-omega}
  For $u \in \Omega$, the lift $\varphi_u$ is an increasing homeomorphism of $\R$, smooth with
  smooth inverse, and both $\varphi_u$ and $\varphi_u^{-1}$ commute with the translation by $1$.
\end{proposition}

\begin{proof}
  \leanok
  $\varphi_u' = 1 + u' > 0$ gives strict monotonicity. Since $u$ is bounded, $\varphi_u(x) \to
  \pm\infty$ as $x \to \pm\infty$, so $\varphi_u$ is surjective by the intermediate value
  theorem. The inverse is smooth by the one-variable inverse function theorem, its derivative
  being $1/\varphi_u' \circ \varphi_u^{-1}$.
\end{proof}

\begin{theorem}[$\Omega$ is closed under the group operations]
  \label{thm:ex-omega-group}
  \lean{LCManifolds.DiffCircle.inv, LCManifolds.DiffCircle.lift_inv, LCManifolds.DiffCircle.inv_mem, LCManifolds.DiffCircle.comp, LCManifolds.DiffCircle.lift_comp, LCManifolds.DiffCircle.comp_mem}
  \leanok
  \uses{prop:ex-lift-diffeo}
  Let $u, w \in \Omega$. Then $\varphi_u^{-1} = \id + v$ with $v \in \Omega$, and
  \[
    \varphi_u \circ \varphi_w = \id + \bigl(w + u \circ \varphi_w\bigr), \qquad
    w + u \circ \varphi_w \in \Omega .
  \]
\end{theorem}

\begin{proof}
  \leanok
  The function $v = \varphi_u^{-1} - \id$ is smooth and periodic by
  Proposition~\ref{prop:ex-lift-diffeo}, and $(\varphi_u^{-1})'(y) =
  1/\varphi_u'(\varphi_u^{-1}y) > 0$. For the composition, $w + u \circ \varphi_w$ is smooth and
  periodic since $\varphi_w$ commutes with the translation by $1$, and by the chain rule
  $1 + (w + u \circ \varphi_w)' = (1 + u' \circ \varphi_w)(1 + w') > 0$.
\end{proof}

Conversely, $\Omega$ captures all the diffeomorphisms of $\R$ that commute with the
translation by $1$, and positivity of the derivative does not need to be assumed.

\begin{theorem}[Characterization of $\Omega$]
  \label{thm:ex-omega-characterization}
  \lean{LCManifolds.DiffCircle.lift_injective, LCManifolds.DiffCircle.mem_domain_iff, LCManifolds.DiffCircle.ofLift, LCManifolds.DiffCircle.existsUnique_mem_domain_lift_eq}
  \leanok
  \uses{prop:ex-lift-diffeo}
  For $u \in C^\infty_{per}$, $u \in \Omega$ if and only if $\varphi_u$ has a differentiable
  left inverse. Consequently, if $\varphi \colon \R \to \R$ is smooth, has a differentiable left
  inverse and satisfies $\varphi(x + 1) = \varphi(x) + 1$, then $\varphi = \id + u$ for a unique
  $u \in \Omega$.
\end{theorem}

\begin{proof}
  \leanok
  A continuous injective map of $\R$ is strictly monotone, and it cannot be decreasing since it
  commutes with the translation by $1$; so $\varphi_u' \ge 0$, and the chain rule
  $\psi'(\varphi_u x)\,\varphi_u'(x) = 1$ excludes zeros. For the second claim take
  $u = \varphi - \id$, which is smooth and periodic; uniqueness is the injectivity of
  $u \mapsto \id + u$.
\end{proof}

So $\Omega$ is exactly the group of diffeomorphisms of $\R$ commuting with the translation by $1$,
that is (by Remark~\ref{rem:ex-covering}), the universal cover of $\Diff^+(S^1)$. Lifts of
orientation reversing diffeomorphisms satisfy $\varphi(x + 1) = \varphi(x) - 1$ instead, and
they are excluded by the hypothesis.

\begin{remark}[What is not proved]
  Theorem~\ref{thm:ex-omega-group} shows that $\Omega$ is closed under composition and
  inversion, so it is a group, and Proposition~\ref{prop:ex-omega-open} shows that it is a
  Fréchet manifold. The smoothness of the maps $\Omega \times \Omega \to \Omega$,
  $(u, w) \mapsto w + u \circ \varphi_w$, and $\Omega \to \Omega$, $u \mapsto v$, is
  \emph{not} proved here; this is what would make $\Omega$ a Lie group, and it belongs to the
  Lie group phase of the project. We remark that no general inverse function theorem is
  available in Fréchet spaces, so the smoothness of the inversion has to be obtained by a
  direct argument.
\end{remark}
